% Deck figure n08: a rotation leaves some parities unchanged (the idea behind README Section 3.2).
% A three-qubit example, bit vectors (x | z): I = (0|0), X = (1|0), Z = (0|1), Y = (1|1).
%   P      = Z I I = 000 | 100   a term
%   Q      = X I I = 100 | 000   the generator of an X rotation on qubit 1
%   P + Q  = Y I I = 100 | 100   the new term that the rotation creates (bitwise sum mod 2)
% The row r is the parity of the three z bits: r(P) = 1, r(Q) = 0, r(P + Q) = 1.
% Because r is linear, r(P + Q) = r(P) + r(Q): a row with r(Q) = 0 has the same value on the new term.
% The example is the one of design_sectors.tex (P1, Q1 and P2 there). No measured numbers in this figure.
\begin{tikzpicture}[x=1cm, y=1cm, font=\rmfamily\footnotesize]
  \tikzset{bc/.style={draw=black!70, line width=0.35pt, anchor=north west, inner sep=0pt, minimum height=0.68cm,
                      text height=0.38cm, text depth=0.12cm, font=\rmfamily\footnotesize},
           bh/.style={bc, fill=black!10, minimum height=0.48cm, text height=0.29cm, text depth=0.09cm},
           bq/.style={bh, minimum height=0.42cm, text height=0.25cm, text depth=0.07cm},
           bm/.style={draw=black!70, line width=0.35pt, anchor=north west, inner sep=0pt, minimum height=0.9cm,
                      fill=black!10, align=center, font=\rmfamily\footnotesize}}
  \node[layer, anchor=north west, minimum width=13.9cm, minimum height=7.05cm] at (0,0) {};
  \node[lttl, anchor=west] at (0.25,-0.48) {A Rotation Creates the Term {\boldmath $P \oplus Q$}};
  % ---- header ----
  \node[bm, minimum width=5.9cm] at (0.75,-1.0) {Pauli String};
  \node[bh, minimum width=1.65cm] at (6.65,-1.0) {Qubit};
  \node[bh, minimum width=1.65cm] at (8.3,-1.0) {x Bits};
  \node[bh, minimum width=1.65cm] at (9.95,-1.0) {z Bits};
  \node[bm, fill=rowc, minimum width=2.05cm] at (11.6,-1.0) {Row $r$ = Parity\\[-1.5pt]of the z Bits};
  \foreach \x in {6.65, 8.3, 9.95}
    \foreach \i in {1,2,3}
      \node[bq, minimum width=0.55cm] at ({\x+0.55*(\i-1)},-1.48) {\i};
  % ---- the term, the generator, the new term ----
  \foreach \y/\n/\d/\fa/\la/\lb/\lc/\xa/\za/\rv in {
      -1.9/{$P$}/{Term}/white/Z/I/I/0/1/{$r(P) = 1$},
      -2.58/{$Q$}/{Generator of an X Rotation on Qubit 1}/genc/X/I/I/1/0/{$r(Q) = 0$},
      -3.68/{$P \oplus Q$}/{New Term}/white/Y/I/I/1/1/{$r(P \oplus Q) = 1$}}
  {
    \node[bc, fill=\fa, minimum width=1.15cm] at (0.75,\y) {\n};
    \node[bc, fill=\fa, minimum width=4.75cm] at (1.9,\y) {\d};
    \foreach \i/\t in {0/\la, 1/\lb, 2/\lc}
      \node[bc, fill=\fa, minimum width=0.55cm] at ({6.65+0.55*\i},\y) {\t};
    \foreach \i/\t in {0/\xa, 1/0, 2/0}
      \node[bc, fill=white, minimum width=0.55cm] at ({8.3+0.55*\i},\y) {\t};
    \foreach \i/\t in {0/\za, 1/0, 2/0}
      \node[bc, fill=rowc!35, minimum width=0.55cm] at ({9.95+0.55*\i},\y) {\t};
    \node[bc, fill=rowc, minimum width=2.05cm] at (11.6,\y) {\rv};
  }
  \node[font=\rmfamily\small] at (0.5,-2.92) {$\oplus$};
  \node[font=\rmfamily\small] at (0.5,-4.02) {$=$};
  \draw[line width=0.7pt] (0.3,-3.47) -- (13.65,-3.47);
  % ---- what changed and what did not ----
  \draw[draw=dred, line width=1.0pt] (8.3,-3.68) rectangle (8.85,-4.36);
  \draw[sarr] (8.575,-4.97) -- (8.575,-4.43);
  \node[note, anchor=base east] at (9.25,-5.25) {Only This x Bit Changed};
  \draw[line width=0.6pt] (9.95,-4.43) -- (9.95,-4.53) -- (11.6,-4.53) -- (11.6,-4.43);
  \draw[sarr] (10.775,-4.97) -- (10.775,-4.57);
  \node[note, anchor=base] at (10.775,-5.25) {z Bits Unchanged};
  % ---- key statement ----
  \node[draw=black!70, line width=0.45pt, fill=white, anchor=north west, minimum width=13.35cm, minimum height=0.9cm] at (0.3,-5.85) {};
  \node[key] at (6.975,-6.3) {If {\boldmath $r(Q) = 0$} Then {\boldmath $r(P \oplus Q) = r(P)$}};
\end{tikzpicture}
